% Source: Solutions/Aaditya S Kannothe-PS4-Q5-8.tex (typed LaTeX), Problem 7
% (student macros expanded: \Rw = \R^\omega, \Bdd = \mathcal{B}, \Unb = \mathcal{U})
\submittedby{Aaditya S Kannothe}

\textbf{Proof.}

The uniform metric is
\[ \bar{\rho}(x,y)=\sup_n\min\{|x_n-y_n|,1\}. \]

Define
\[ \mathcal{B}=\left\{x\in\R^{\omega}:\sup_n|x_n|<\infty\right\} \]
and
\[ \mathcal{U}=\R^{\omega}\setminus\mathcal{B}. \]

We show that both sets are open.

First, observe that
\[ \bar{\rho}(x,y)<1\implies |x_n-y_n|<1\quad\forall n. \]

\textbf{Claim 1: $\mathcal{B}$ is open.}

Let $x\in\mathcal{B}$. Then $|x_n|\leq M$ for some $M>0$ and all $n$.

If $\bar{\rho}(x,y)<1$, then
\[ |y_n|\leq |x_n|+|y_n-x_n|<M+1. \]
Hence $y\in\mathcal{B}$, proving that $\mathcal{B}$ is open.

\textbf{Claim 2: $\mathcal{U}$ is open.}

Let $x\in\mathcal{U}$ and suppose $\bar{\rho}(x,y)<1$.

If $y$ were bounded, say $|y_n|\leq M$, then
\[ |x_n|\leq |y_n|+|x_n-y_n|<M+1. \]
This contradicts the fact that $x$ is unbounded.

Thus $y\in\mathcal{U}$, proving that $\mathcal{U}$ is open.

Finally, $\mathcal{B}$ and $\mathcal{U}$ are disjoint, non-empty, and
\[ \mathcal{B}\cup\mathcal{U}=\R^{\omega}. \]

Therefore, they form a separation of $\R^{\omega}$. Hence,
\[ \boxed{\R^{\omega}\text{ is not connected.}} \]
\hfill $\square$
